\documentclass[11pt]{article}
\input{style-pc}

\begin{document}

\entetesujet{Sujet bac 2010 -- Physique-Chimie -- Série D}

% ============================ CHIMIE ============================
\matiere{CHIMIE}{8 points}

\exo{1}{}
On rappelle que les niveaux d'énergie quantifiés de l'atome d'hydrogène sont donnés par la relation :
\[ E_n = -\frac{E_0}{n^2} \quad \text{avec } E_0 = 13{,}6 \text{ eV et } n \text{ étant un entier positif.} \]

\begin{enumerate}
  \item \begin{enumerate}[label=\alph*.]
    \item Déterminer, en joules, l'énergie qu'il faut fournir à un atome d'hydrogène pour permettre son passage de l'état fondamental au premier état excité.
    \item Que se passe-t-il si l'atome d'hydrogène dans son état fondamental reçoit :
    \begin{itemize}
      \item un photon d'énergie $W = 1{,}83\cdot10^{-18}$ J ?
      \item un électron d'énergie cinétique $E_c = 1{,}83\cdot10^{-18}$ J ?
    \end{itemize}
  \end{enumerate}
  \item Définir et calculer l'énergie d'ionisation de l'atome d'hydrogène.
  \item On considère la série de Lyman.
  \begin{enumerate}[label=\alph*.]
    \item Qu'appelle-t-on série de raies ?
    \item L'analyse spectroscopique permet de déceler la radiation de fréquence $\nu = 3{,}8\cdot10^{15}$ Hz. À quelle transition correspond-elle ?
  \end{enumerate}
\end{enumerate}
\noindent On donne : $h = 6{,}62\cdot10^{-34}$ J$\cdot$s ; $c = 3\cdot10^8$ m$\cdot$s$^{-1}$ ; $1$ eV $= 1{,}6\cdot10^{-19}$ J.

\exo{2}{}
On dose une eau de Javel à usage domestique. Pour cela, on fait réagir 20 mL de cette eau de Javel diluée contenant des ions hypochlorite ($\ch{ClO^-}$) dans un excès d'ions iodures $\ch{I^-}$. On acidifie le milieu.

\begin{enumerate}
  \item Sachant que les couples redox en présence sont : $\ch{ClO^-/Cl^-}$ et $\ch{I_2/I^-}$, écrire l'équation bilan de la réaction redox.
  \item On dose les molécules de diiode $\ch{I_2}$ formées par une solution de thiosulfate ($\ch{S_2O_3^{2-}}$) de concentration molaire $0{,}100$ mol$\cdot$L$^{-1}$. L'équivalence est atteinte pour $15{,}2$ mL de solution de thiosulfate versés.
  \begin{enumerate}[label=\alph*.]
    \item Écrire l'équation bilan de la réaction de dosage. On donne le couple redox $(\ch{S_4O_6^{2-}/S_2O_3^{2-}})$.
    \item Calculer la concentration molaire de l'eau de Javel en ions $\ch{ClO^-}$.
  \end{enumerate}
\end{enumerate}
\noindent N.B : les molécules de diiode ($\ch{I_2}$) sont à l'état liquide à la température de l'expérience.

% ============================ PHYSIQUE ============================
\matiere{PHYSIQUE}{12 points}

\exo{1}{}
Un proton animé d'une vitesse $\vec{V_0}$ entre dans un champ électrostatique uniforme $\vec{E}$ créé entre deux plaques $A$ et $B$, par un point $O$ situé à égale distance des plaques. La différence de potentiel entre les plaques est $U = 400$ V. On néglige le poids du proton.

\begin{center}
\begin{tikzpicture}[>=Stealth,scale=1]
  % axes
  \draw[->] (0,-0.2)--(0,2) node[above]{$y$};
  \draw[->] (0,0)--(8.4,0) node[right]{$x$};
  % plaques
  \draw[thick] (0.3,1)--(4.3,1) node[right]{$A$};
  \draw[thick] (0.3,-1)--(4.3,-1) node[right]{$B$};
  % O et V0
  \node[below left] at (0,0){$O$};
  \draw[->,thick] (0,0)--(1.2,0) node[above]{$\vec{V_0}$};
  % champ E
  \draw[->] (2.6,0.4)--(2.6,-0.4) node[right]{$\vec{E}$};
  % ecran
  \draw[thick] (8,-1.6)--(8,1.6);
  % d
  \draw[<->,dashed] (-0.3,-1)--(-0.3,1) node[midway,left]{$d$};
  % L
  \draw[<->,dashed] (0.3,-1.45)--(4.3,-1.45) node[midway,below]{$L$};
  % D
  \draw[<->,dashed] (2.6,-0.55)--(8,-0.55) node[midway,below]{$D$};
\end{tikzpicture}
\end{center}

\begin{enumerate}
  \item Sur un schéma clair :
  \begin{enumerate}[label=\alph*.]
    \item Indiquer les signes des plaques. Justifier.
    \item Représenter la force électrostatique $\vec{F}$ qui s'exerce sur le proton dans le champ électrostatique.
  \end{enumerate}
  \item \begin{enumerate}[label=\alph*.]
    \item Établir les équations horaires du mouvement du proton dans le repère $(O, x, y)$.
    \item En déduire l'équation de la trajectoire du proton à l'intérieur des plaques.
  \end{enumerate}
  \item Le proton sort du champ par le point $S$ d'ordonnée $y_S = -0{,}96$ mm.
  \begin{enumerate}[label=\alph*.]
    \item Déterminer $V_0$.
    \item Quelle est la nature du mouvement à l'extérieur des plaques ?
  \end{enumerate}
  \item On place un écran vertical à une distance $D = 30$ cm du milieu des plaques. Déterminer les coordonnées du point d'impact $M$ du proton sur l'écran.
\end{enumerate}
\noindent On donne : $d = 20$ cm ; $L = 10$ cm ; masse du proton : $m_p = 1{,}67\cdot10^{-27}$ kg ; charge du proton : $q_p = +e$ ; $e = 1{,}6\cdot10^{-19}$ C.

\exo{2}{}
\begin{minipage}{0.66\linewidth}
Un pendule simple est constitué d'un fil inextensible de masse négligeable et de longueur $l$ et d'un solide ponctuel de masse $m = 100$ g. L'extrémité libre du fil est fixée à un point $O$ d'un support.

Le pendule effectue des oscillations de faible amplitude, de période $T = 2$ s autour de l'axe horizontal $(\Delta)$ passant par le point $O$.
\end{minipage}\hfill
\begin{minipage}{0.3\linewidth}
\centering
\begin{tikzpicture}[>=Stealth,scale=1]
  \draw[thick] (0,0)--(0,-3.2);
  \node[above] at (0,0){$O$};
  \node[right] at (0.05,-0.05){$(\Delta)$};
  \draw (-0.12,-0.12)--(0.12,0.12); \draw (-0.12,0.12)--(0.12,-0.12);
  \fill (0,-3.2) circle (2pt) node[below]{$A$};
\end{tikzpicture}
\end{minipage}

\begin{enumerate}
  \item Calculer :
  \begin{enumerate}[label=\alph*.]
    \item La longueur du fil.
    \item L'incertitude absolue sur la mesure de la longueur sachant que $g = 9{,}80 \pm 0{,}01$ m/s$^2$ et que la période a été mesurée à $0{,}02$ s près.
  \end{enumerate}
  \item On écarte le pendule de sa position d'équilibre d'un angle de $60^\circ$ et on l'abandonne sans vitesse initiale. Déterminer :
  \begin{enumerate}[label=\alph*.]
    \item En appliquant le théorème de l'énergie cinétique, la vitesse linéaire du solide à son passage par la position $\theta = 45^\circ$.
    \item La tension du fil à cette position.
  \end{enumerate}
  \item Calculer l'énergie cinétique du solide lorsqu'il passe par la position verticale.
\end{enumerate}
\noindent N.B : pour les questions 2 et 3, on prendra $g = 9{,}8$ m$\cdot$s$^{-2}$.

\end{document}
